\section{Reutilización de código y DRY}

En escenarios de Vibe Coding, el modelo está limitado por la ventana de contexto, por lo que tiende a modificar el código habiendo visto solo una parte de él, lo que provoca que las funciones de utilidad (Utils) acumulen una gran cantidad de código redundante y duplicado.

Superpowers insiste una y otra vez, tanto en \texttt{Writing Plans} como en la revisión de código:

\begin{itemize}
    \item \textbf{DRY} (Don't Repeat Yourself): eliminar el código duplicado.
    \item \textbf{YAGNI} (You Aren't Gonna Need It): eliminar de manera decidida el diseño excesivo y la abstracción excesiva.
\end{itemize}

Además, Superpowers ofrece el Skill \texttt{Writing Skills}, que insiste en depurar el proceso de resolución de problemas hasta convertirlo en un \textbf{patrón reutilizable}, en lugar de un script de un solo uso.

\subsection{Resumen de la sección}

DRY y YAGNI son los dos principios que atraviesan todo Superpowers: el primero evita la redundancia y el segundo evita la sobreingeniería; juntos constituyen el equilibrio dinámico de la calidad del código.

